\documentclass[../../main2.tex]{subfiles}

\begin{document}

	% ======================================================================
	% v2 新 Part III：法律的角色从「四个通道之一」升格为
	% 「因果贡献分解方法最早的成熟应用实例」——即方法的证明，而非又一个通道。
	% 与 Part IV 的 chapter12_law.tex（法律作为四大齿轮之一）刻意区分：
	%   本章看「方法怎么被逼出来」，Part IV 看「它在系统里咬合在哪」。
	% 正文由后续写作阶段填充；灰色为骨架追问链与待写标记。
	% ======================================================================

	\chapter{法庭上的时光机：当因果必须给出一个数字}
	\label{ch:law-method}

	\begin{motivation}
		Part II 给了我们一个工具：把「A 和 B 共同导致 D，各占多少」写成一个可比的量——\\
		$C_A = P(D \mid A) - P(D \mid \neg A)$，也就是「无你检验」：把 A 拿掉，D 的概率掉多少，掉的那块就是 A 的功劳。\\
		\textbf{核心动机}：你可能会想，这是书斋里的数学玩具。这一章要证明它不是——\\
		有一群人被迫每天使用这个工具，而且他们不能给模糊答案：必须说清「70\% 是他的责任」。\\
		这群人叫法官。法律，是因果贡献分解方法最早的成熟应用现场。
	\end{motivation}

	\begin{logicmap}
	\begin{tikzpicture}[node distance=0.9cm, every node/.style={draw, rectangle, rounded corners, align=center}]
	\node (q) {「各占多少」只是数学吗？};
	\node (c) [below=of q] {法庭：必须给确定数字（§法1）};
	\node (m) [below=of c] {倒逼出精细的因果分解（§法2）};
	\node (p) [below=of m] {与 Part II 工具同构（§法3）};
	\node (n) [below=of p] {方法被真实世界验过（抛给 Part IV）};
	\draw[-{Stealth}] (q) -- (c);
	\draw[-{Stealth}] (c) -- (m);
	\draw[-{Stealth}] (m) -- (p);
	\draw[-{Stealth}] (p) -- (n);
	\end{tikzpicture}
	\end{logicmap}

	% ==================================================================
	\section{当「差不多」不行的时候}
	\label{sec:court-must-number}

	\textcolor{gray}{\textit{【骨架 · 追问链（终版）】为什么法庭是检验因果分解的天然实验室？→ 因为纠纷必须了结：不能停留在「有点像他的错」，必须说出「70\% 他的责任」。这个「必须给数字」的压力，是科学界直到近代才有的，法律界却扛了几千年。}}

	\textcolor{gray}{（正文待写：用一句日常钩子开场——「你有没有想过，为什么『不得杀人』四个字人人会背，真要定责时却要几百页卷宗？」；对比科学里的 p<0.05（模糊门槛）与法庭里的责任比例（确定数字）；点明本章主旨：本书的工具不是发明，是被现实逼出来的。）}

	\paragraph*{逻辑衔接}
	\textcolor{gray}{（待写）「必须给数字」催生了什么？下一节：一套关于「怎么分锅」的精细方法。}

	% ==================================================================
	\section{把责任拆成可比较的量}
	\label{sec:court-decompose}

	\textcolor{gray}{\textit{【骨架 · 追问链（终版）】法学从「故意 / 过失」二分类，推进到「责任比例」——这正是条件概率贡献分解的现实版本。法律不必懂 $P(D\mid A)$，但它做的事和它同构：把多个致害因素各自「拿掉一次」，看损害概率掉多少。}}

	\textcolor{gray}{（正文待写：用真实判例（如多车连环相撞、产品缺陷致人损害）说明「多因搅在一起」如何被拆解；「举证责任」「优势证据」的实质就是粗粒化的概率差门槛；诚实标注：显式概率化定责在实务中罕见，但逻辑内核一致——这是本书不回避的「理想化模型」。）}

	\begin{questionbox}
		\textcolor{gray}{（待写：反驳位——「法官定责是拍脑袋，哪真算概率？」回答：即使不写公式，判决的推理结构仍是「若无被告行为，损害是否仍会发生」，这正是 but-for test / 无你检验。）}
	\end{questionbox}

	\paragraph*{逻辑衔接}
	\textcolor{gray}{（待写）法庭的方法，和 Part II 的工具是不是同一个东西？下一节做同构对照——这是本章的高潮。}

	% ==================================================================
	\section{与「无你检验」的同构}
	\label{sec:court-isomorphism}

	\textcolor{gray}{\textit{【骨架 · 追问链（终版）】Part II 的公式 $C_A = P(D\mid A)-P(D\mid\neg A)$ = 法律定责公式。一旦你接受「责任 $\propto P(\text{损害}\mid\text{行为})-P(\text{损害}\mid\neg\text{行为})$」，你就得到了社会科学因果推断的统一语言。法律在这里不是被套用的例子，是方法的发源地。}}

	\textcolor{gray}{（正文待写：把 but-for test 逐符号对应到 $C_A$；说明「相当因果关系」「过错相抵」如何翻译为贡献的加减；这是全书两半（工具 / 应用）的真正接合点——而它发生在法庭，不在实验室。）}

	\paragraph*{逻辑衔接}
	\textcolor{gray}{（待写）方法被真实世界验过了。但法庭只处理「已经发生的纠纷」；我们要预测的是「还没发生的社会走向」。下一章（Part IV）把同一个工具，装到社会运转的四个齿轮上。}

	% ==================================================================
	\section{本章小结与衔接}
	\label{sec:court-summary}

	\textcolor{gray}{（正文待写：收束 Part III——一句话：因果贡献分解不是本书的发明，是人类在法庭上被迫练了几千年的基本功；本书只是把它从「判案」推广到「预测」。）}

	\paragraph*{逻辑衔接}
	\textcolor{gray}{（待写）带着这套已被验证的工具进入 Part IV：社会的因果，具体长什么样？}

\end{document}
